\chapter{Algebraic Data Types: Making Illegal States Unrepresentable}
\label{ch05:adts}

Chapters~3 and~4 tracked \emph{who owns} a block of memory and \emph{how long} a borrow may point into it. Ownership and lifetimes rule out dangling pointers, double frees, and iterator invalidation, yet a memory-safe program can still wire money twice, settle a canceled trade, or panic at three in the morning when an unexpected \texttt{null} reaches a report generator. Those bugs live one level up: in the gap between the bit patterns a data structure \emph{can} hold in RAM and the states that actually make sense in the problem domain.

If you grew up in Java~6 and C++, your reflexes for modeling domain data were shaped by mutable classes full of nullable references and boolean flags, plus open \texttt{extends} hierarchies on the heap. Rust drops both \texttt{extends} implementation inheritance and implicit \texttt{null}. In their place, it gives you \emph{algebraic data types} (\texttt{struct} and data-carrying \texttt{enum}) paired with exhaustive pattern matching (\texttt{match}). This chapter counts how many nonsensical states each design allows, and then inspects the exact bytes the compiler emits for the safer design.

We will study Rust's data model from two angles: first from above, using simple counting algebra to measure how tightly a type fits its domain, and then down to the metal, counting exact bytes, alignment padding, discriminant tags, and bit-pattern niches on a 64-bit machine.

\chapterline

\section{Why Rust Dropped \texttt{extends} and \texttt{null}: The Algebra of Cardinality}
\label{ch05:cardinality}

Every type defines the set of values a variable of that type can hold at runtime. Before asking what methods a type offers, a Rust engineer asks a simpler counting question: \emph{how many distinct values inhabit this type?}

\begin{definition}[Cardinality of a type]
\label{ch05:def-cardinality}
For any type \texttt{T}, its \emph{cardinality}, written $\abs{\texttt{T}}$, is the number of distinct valid values of type \texttt{T} (ignoring finite RAM limits on heap containers).
\end{definition}

Count the inhabitants of Rust's smallest built-in primitives:
\begin{itemize}
  \item \textbf{Booleans and fixed-width integers.} A \texttt{bool} has $\abs{\texttt{bool}} = 2$ (\texttt{false}, \texttt{true}); a \texttt{u8} has $\abs{\texttt{u8}} = 2^8 = 256$; a \texttt{u64} has $\abs{\texttt{u64}} = 2^{64}$.
  \item \textbf{The unit type \texttt{()}: $\abs{\texttt{()}} = 1$.} Its sole value is also written \texttt{()}. Where C and Java use \texttt{void} for ``returns nothing,'' creating an awkward split between \texttt{void} methods and generic functions returning \texttt{T}, Rust functions without a return arrow return \texttt{()}. Because only one value exists, carrying a \texttt{()} conveys zero runtime bits of information, so \texttt{core::mem::size\_of::<()>() == 0}.
  \item \textbf{The never type \texttt{!}: $\abs{\texttt{!}} = 0$.} Rust also has a type with \emph{zero} inhabitants, written \texttt{!}. Expressions that diverge before producing a value --- \texttt{return}, \texttt{break}, \texttt{continue}, \texttt{panic!(...)}, \texttt{std::process::abort()}, and an infinite \texttt{loop \{\}} --- all have type \texttt{!}.
\end{itemize}

A zero-inhabitant type is useful every time a \texttt{match} mixes value-producing arms with early exits:

\begin{lstlisting}[language=Rust]
fn unwrap_or_abort<T>(opt: Option<T>) -> T {
    match opt {
        Some(value) => value,            // produces T
        None => std::process::abort(),   // produces !, which coerces to any T
    }
}
\end{lstlisting}

Every arm of a \texttt{match} must agree on one result type \texttt{T}. Because \texttt{abort()} has type \texttt{!} ($\abs{\texttt{!}} = 0$), control never reaches the end of the second arm. Just as the empty set $\emptyset$ is a subset of every set $S$, the uninhabited type \texttt{!} coerces soundly to any type \texttt{T}.

\subsection{Product Types: \texttt{struct} Multiplies Possibilities}
\label{ch05:product-types}

Bundling fields \texttt{a: A} and \texttt{b: B} into a record lets every value of \texttt{A} pair independently with every value of \texttt{B}. That set is the Cartesian product $A \times B$, so cardinalities multiply:
\[
  \abs{\texttt{struct \{ a: A, b: B \}}} = \abs{A \times B} = \abs{A} \cdot \abs{B}.
\]
Because their cardinalities multiply, structs and tuples are \emph{product types}. Rust provides three forms:
\begin{enumerate}
  \item \textbf{Named-field structs}, such as \texttt{struct Point \{ x: i32, y: i32 \}}, with cardinality $2^{32} \cdot 2^{32} = 2^{64}$.
  \item \textbf{Tuple structs}, such as \texttt{pub struct BasisPoints(pub u16);}, with cardinality $65{,}536$. Wrapping a primitive in a single-field tuple struct --- the \emph{newtype} idiom --- creates a distinct compile-time type with \emph{zero} heap boxing or object header: in machine code it is an unboxed 2-byte \texttt{u16} in a register, yet the compiler rejects passing a bare \texttt{u16} where \texttt{BasisPoints} are expected.
  \item \textbf{Unit structs}, such as \texttt{pub struct Usd;}, an empty product ($\abs{\texttt{Usd}} = 1$, zero runtime bytes) used as type-level markers in Chapter~6.
\end{enumerate}

Multiplication is how flag-and-nullable classes end up allowing far more states than the domain has. Here is a typical settlement record from a Java~6 back-office service:

\begin{lstlisting}[language=Java]
// Java 6: A four-field product of booleans and nullable references
public final class SettlementRecord {
    private final boolean isSettled;
    private final boolean isFailed;
    private final Long settledAtEpochMs; // null until settled
    private final Integer failureCode;   // null unless failed

    public SettlementRecord(
            boolean isSettled, boolean isFailed,
            Long settledAtEpochMs, Integer failureCode) {
        this.isSettled = isSettled;
        this.isFailed = isFailed;
        this.settledAtEpochMs = settledAtEpochMs;
        this.failureCode = failureCode;
    }
}
\end{lstlisting}

Count the domain states first. A transfer is in \emph{one} of three mutually exclusive stages: \textbf{Pending} ($1$ state), \textbf{Settled} at a 64-bit timestamp ($2^{64}$ states), or \textbf{Failed} with a 32-bit code ($2^{32}$ states) --- a total of $1 + 2^{64} + 2^{32}$ meaningful states.

Now count what \texttt{SettlementRecord} actually represents. Because every Java object reference includes \texttt{null}, $\abs{\texttt{Long}} = 1 + 2^{64}$ and $\abs{\texttt{Integer}} = 1 + 2^{32}$. Multiplying the four fields yields:
\[
  2 \cdot 2 \cdot (1 + 2^{64}) \cdot (1 + 2^{32}) \;=\; 4 \cdot (1 + 2^{32} + 2^{64} + 2^{96}) \;\approx\; 2^{98}.
\]
Fewer than \emph{one in ten billion} representable configurations of \texttt{SettlementRecord} are valid. All the rest --- setting both flags to \texttt{true}, marking \texttt{isSettled} while leaving \texttt{settledAtEpochMs} as \texttt{null} (which throws on unboxing), or leaving a stale \texttt{failureCode} on a pending transfer --- force defensive \texttt{if} checks into every method.

\subsection{Sum Types: \texttt{enum} Adds Possibilities}
\label{ch05:sum-types}

Look back at the true domain count: $1 + 2^{64} + 2^{32}$. Those are \emph{plus} signs, not multiplication signs. What the domain describes is a disjoint union where each alternative carries \emph{only} its own payload. In Rust, an \texttt{enum} defines a closed \emph{sum type} whose variants may be unit-like, tuple-like, or struct-like:

\begin{lstlisting}[language=Rust]
pub enum SettlementStatus {
    Pending,
    Settled { settled_at_epoch_ms: u64 },
    Failed { failure_code: u32 },
}
\end{lstlisting}

Its cardinality is the exact sum $\abs{\texttt{()}} + \abs{\texttt{u64}} + \abs{\texttt{u32}} = 1 + 2^{64} + 2^{32}$. None of the illegal states can be represented: a value holds only one tag at a time, \texttt{Settled} cannot be built without a \texttt{u64} timestamp, and \texttt{failure\_code} does not even exist on \texttt{Settled}. Figure~\ref{ch05:fig-product-vs-sum} compares the Cartesian grid of a flag product against the disjoint union of a sum type.

\begin{figure}[H]
\centering
\begin{tikzpicture}[font=\small,>=Latex]
  % Left panel: Product Type (x = 0 .. 5.8)
  \begin{scope}[shift={(0,0)}]
    \node[font=\bfseries,sourceblue] at (2.95,4.15) {Product: $|A \times B| = |A| \cdot |B|$};
    \node[font=\footnotesize,black!75] at (2.95,3.72) {Flags \texttt{(settled, failed)}: $2 \times 2 = 4$ cells};

    \node[font=\scriptsize\ttfamily] at (2.05,3.15) {failed=false};
    \node[font=\scriptsize\ttfamily] at (4.55,3.15) {failed=true};
    \node[font=\scriptsize\ttfamily,anchor=east] at (0.85,2.25) {settled=false};
    \node[font=\scriptsize\ttfamily,anchor=east] at (0.85,1.05) {settled=true};

    \draw[thick,sourceblue!70,fill=sourceblue!7,rounded corners=3pt] (0.95,1.70) rectangle (3.15,2.80);
    \node[font=\footnotesize\bfseries] at (2.05,2.42) {Pending};
    \node[font=\scriptsize,black!75] at (2.05,2.00) {valid};

    \draw[thick,sourceblue!70,fill=sourceblue!7,rounded corners=3pt] (3.35,1.70) rectangle (5.75,2.80);
    \node[font=\footnotesize\bfseries] at (4.55,2.42) {Failed};
    \node[font=\scriptsize\ttfamily,black!75] at (4.55,2.00) {code != null};

    \draw[thick,sourceblue!70,fill=sourceblue!7,rounded corners=3pt] (0.95,0.50) rectangle (3.15,1.60);
    \node[font=\footnotesize\bfseries] at (2.05,1.22) {Settled};
    \node[font=\scriptsize\ttfamily,black!75] at (2.05,0.80) {time != null};

    \draw[very thick,trapred,fill=trapred!12,rounded corners=3pt] (3.35,0.50) rectangle (5.75,1.60);
    \node[font=\footnotesize\bfseries,trapred] at (4.55,1.22) {Both true};
    \node[font=\scriptsize,trapred] at (4.55,0.80) {illegal state!};
  \end{scope}

  % Divider
  \draw[dashed,gray!60,thick] (6.30,0.40) -- (6.30,4.35);

  % Right panel: Sum Type (x = 6.80 .. 13.30)
  \begin{scope}[shift={(6.80,0)}]
    \node[font=\bfseries,targetgreen!85!black] at (3.25,4.15) {Sum: $|A + B + C| = |A| + |B| + |C|$};
    \node[font=\footnotesize,black!75] at (3.25,3.72) {\texttt{enum SettlementStatus}: disjoint variants};

    \draw[thick,targetgreen!80!black,fill=targetgreen!8,rounded corners=4pt] (0.10,2.60) rectangle (6.40,3.30);
    \node[font=\footnotesize\ttfamily,anchor=west] at (0.30,2.95) {Pending};
    \node[font=\scriptsize,targetgreen!75!black,anchor=east] at (6.20,2.95) {$1$ value};

    \draw[thick,targetgreen!80!black,fill=targetgreen!8,rounded corners=4pt] (0.10,1.55) rectangle (6.40,2.35);
    \node[font=\footnotesize\ttfamily,anchor=west] at (0.30,1.95) {Settled \{ settled\_at: u64 \}};
    \node[font=\scriptsize,targetgreen!75!black,anchor=east] at (6.20,1.95) {$2^{64}$ values};

    \draw[thick,targetgreen!80!black,fill=targetgreen!8,rounded corners=4pt] (0.10,0.50) rectangle (6.40,1.30);
    \node[font=\footnotesize\ttfamily,anchor=west] at (0.30,0.90) {Failed \{ code: u32 \}};
    \node[font=\scriptsize,targetgreen!75!black,anchor=east] at (6.20,0.90) {$2^{32}$ values};
  \end{scope}
\end{tikzpicture}
\caption{Product types multiply fields into a Cartesian grid with illegal cells such as \texttt{(settled=true, failed=true)}, whereas sum types add disjoint variants carrying only their own non-null fields.}
\label{ch05:fig-product-vs-sum}
\end{figure}

\begin{intuitionbox}
Reading $\times$ as \texttt{struct}/tuple, $+$ as \texttt{enum}, $1$ as \texttt{()}, and $0$ as \texttt{!} turns ordinary algebra into type isomorphisms:
\begin{itemize}
  \item $T \times 1 = T$: \texttt{(u64, ())} carries the exact same $2^{64}$ states and 8 bytes as \texttt{u64}.
  \item $T + 0 = T$: \texttt{Result<u64, !>} has $2^{64} + 0 = 2^{64}$ inhabitants because \texttt{Err(!)} is unconstructible.
  \item $T \times 0 = 0$: \texttt{struct Impossible \{ id: u64, proof: ! \}} has $2^{64} \cdot 0 = 0$ inhabitants; nobody can ever construct it without first producing a value of \texttt{!}.
  \item $1 + T$: \texttt{Option<T>} adds \emph{one} explicit absence state (\texttt{None}) only to types that ask for it.
\end{itemize}
\end{intuitionbox}

\subsection{Why Java~6 Subclasses and C \texttt{union}s Fall Short}
\label{ch05:java-c-contrast}

Java~6 and C each offer half of what Rust's \texttt{enum} provides:
\begin{itemize}
  \item \textbf{Java~6 \texttt{enum} + \texttt{abstract class}.} Java's \texttt{enum} constants are fixed singletons allocated at class load time, so \texttt{SETTLED} cannot hold a per-transfer timestamp. Modeling per-case data in Java~6 forces an \texttt{abstract class SettlementStatus} with heap-allocated subclasses (\texttt{PendingStatus}, \texttt{SettledStatus}, \texttt{FailedStatus}). Every field now chases a pointer to a separate GC object with a 12-to-16-byte header, the reference can still be \texttt{null}, and open subclassing across packages defeats compiler-checked exhaustiveness unless you write clunky Visitor boilerplate. (Chapter~10 examines how Kotlin's \texttt{sealed} classes modernize the JVM side of this story.)
  \item \textbf{C's unchecked \texttt{union}.} A C \texttt{struct} pairing an \texttt{enum} tag with a \texttt{union \{ uint64\_t settled\_at; uint32\_t failure\_code; \}} sits flat inline in memory with zero pointer indirection, just like systems code demands:
\end{itemize}

\begin{lstlisting}[language=C]
enum StatusTag { TAG_PENDING, TAG_SETTLED, TAG_FAILED };
struct SettlementC {
    enum StatusTag tag;
    union {
        uint64_t settled_at_epoch_ms;
        uint32_t failure_code;
    } payload;
};
\end{lstlisting}

Reading \texttt{rec.payload.settled\_at\_epoch\_ms} while \texttt{rec.tag} holds \texttt{TAG\_FAILED} compiles in C and silently reinterprets an error code plus four uninitialized bytes as a timestamp. Rust's \texttt{enum} compiles down to that same flat tagged-union layout while making it impossible in safe code to touch a variant's payload without matching its tag first.

\begin{historybox}
Closed sum types with exhaustive pattern matching come from the Edinburgh functional-programming tradition. \textbf{Robin Milner}'s \textbf{ML} (1973) was built as the scripting language of the LCF theorem prover, and the related language \textbf{Hope} (Burstall, MacQueen, and Sannella, 1980) introduced user-defined algebraic data types with pattern matching, which Standard ML then adopted. Mainstream systems languages followed a different path. In 1965 \textbf{Tony Hoare} added a \texttt{null} reference to \textbf{ALGOL~W}, by his own account ``simply because it was so easy to implement,'' and C, C++, and Java inherited the idea. At QCon London in 2009 he called it his ``billion-dollar mistake.'' Rust joins the two traditions: ML-style algebraic types, compiled down to C-style flat memory and bit-pattern niches.
\end{historybox}

\section{Exact Byte Layout: Struct Reordering, Tagged Unions, and Niche Optimization}
\label{ch05:memory-layout}

Every sized type \texttt{T} in Rust has two compile-time constants, which you can read with the functions \texttt{size\_of} and \texttt{align\_of} from \texttt{core::mem}:
\begin{itemize}
  \item \textbf{Alignment (\texttt{align\_of::<T>()}):} a power of two ($1, 2, 4, 8, \dots$) dictating that every valid address of \texttt{T} must be a multiple of that byte count. On 64-bit targets, \texttt{u8}/\texttt{bool} have alignment~$1$, \texttt{u16} has~$2$, \texttt{u32} has~$4$, and \texttt{u64}, \texttt{usize}, \texttt{\&T}, and \texttt{Box<T>} have~$8$.
  \item \textbf{Size (\texttt{size\_of::<T>()}):} the byte stride between elements in \texttt{[T; N]}, always rounded up to a multiple of \texttt{align\_of::<T>()}.
\end{itemize}

\subsection{Struct Field Reordering: \texttt{repr(Rust)} vs.\ \texttt{\#[repr(C)]}}
\label{ch05:struct-reordering}

C and C++ require struct fields to sit in memory in the exact order written in source code. Alternating small and wide fields forces dead alignment padding between them. Compare \texttt{\#[repr(C)]} against Rust's default \texttt{repr(Rust)} on 64-bit hardware:

\begin{lstlisting}[language=Rust]
use core::mem::{align_of, size_of};

#[repr(C)]
pub struct PaddedHeaderC {
    pub version: u8,   // byte 0   (1B) + 7B padding so tx_id lands on byte 8
    pub tx_id: u64,    // bytes 8..=15  (8B)
    pub flags: u8,     // byte 16  (1B) + 7B tail padding so stride % 8 == 0
}

pub struct CompactHeaderRust {
    pub version: u8,
    pub tx_id: u64,
    pub flags: u8,
}

const _: () = assert!(size_of::<PaddedHeaderC>() == 24);
const _: () = assert!(align_of::<PaddedHeaderC>() == 8);
const _: () = assert!(size_of::<CompactHeaderRust>() == 16);
const _: () = assert!(align_of::<CompactHeaderRust>() == 8);
\end{lstlisting}

Because \texttt{size\_of} and \texttt{align\_of} are \texttt{const fn}s, \texttt{const \_: () = assert!(...);} checks layouts at compile time at zero runtime cost. Under \texttt{\#[repr(C)]}, 10 bytes of actual payload expand to 24 bytes (14 bytes of padding) so \texttt{tx\_id} aligns to 8 bytes in every element of \texttt{[PaddedHeaderC; N]}. Under default \texttt{repr(Rust)}, field offsets are intentionally unspecified so \texttt{rustc} can place \texttt{tx\_id: u64} first at bytes \texttt{0..=7}, followed by \texttt{version} at byte~8, \texttt{flags} at byte~9, and 6 tail-padding bytes (\texttt{10..=15}) --- saving 33\% of cache space without forcing you to reorder source fields for the machine. Add \texttt{\#[repr(C)]} only at C FFI or kernel ABI boundaries.

\subsection{How a General \texttt{enum} Sits in Memory: The Tagged Union}
\label{ch05:tagged-union-layout}

Now examine how \texttt{rustc} lays out a multi-variant domain enum:

\begin{lstlisting}[language=Rust]
use core::mem::{align_of, size_of};

#[repr(C, u8)]
pub enum LedgerEvent {
    Heartbeat,                                             // tag = 0
    WirePosted { tx_id: u64, fee_bps: u16 },               // tag = 1
    BatchClosed { batch_id: u64, count: u32, flags: u32 }, // tag = 2
}

const _: () = assert!(size_of::<LedgerEvent>() == 24);
const _: () = assert!(align_of::<LedgerEvent>() == 8);
\end{lstlisting}

Adding \texttt{\#[repr(C, u8)]} pins a 1-byte discriminant tag at byte~0 followed by a C-compatible union of the variant structs (omitting the attribute produces the same 24-byte size here). Figure~\ref{ch05:fig-enum-byte-layout} maps all 24 bytes (\texttt{0..=23}): byte~\texttt{0} holds tag \texttt{0}, \texttt{1}, or \texttt{2}; bytes \texttt{1..=7} pad out to an 8-byte boundary for the \texttt{u64} fields; and bytes \texttt{8..=23} overlay whichever variant is active, sized to the 16-byte \texttt{BatchClosed} payload.

\begin{figure}[H]
\centering
\begin{tikzpicture}[font=\footnotesize,>=Latex]
  \def\xzero{2.55}
  \def\bw{0.48}

  \node[font=\scriptsize\bfseries,black!70,anchor=east] at (\xzero-0.15,3.55) {Byte offset:};
  \node[font=\scriptsize\ttfamily] at ({\xzero+0.5*\bw},3.55) {0};
  \node[font=\scriptsize\ttfamily] at ({\xzero+4.0*\bw},3.55) {1..7};
  \node[font=\scriptsize\ttfamily] at ({\xzero+12.0*\bw},3.55) {8..15};
  \node[font=\scriptsize\ttfamily] at ({\xzero+20.0*\bw},3.55) {16..23};

  % Row 1: Heartbeat
  \node[font=\footnotesize\ttfamily\bfseries,anchor=east] at (\xzero-0.15,2.80) {Heartbeat};
  \draw[thick,sourceblue,fill=sourceblue!18] (\xzero,2.45) rectangle ({\xzero+1*\bw},3.15);
  \node[font=\scriptsize\ttfamily\bfseries,sourceblue!85!black] at ({\xzero+0.5*\bw},2.80) {0};
  \draw[thick,gray!65,fill=gray!12] ({\xzero+1*\bw},2.45) rectangle ({\xzero+24*\bw},3.15);
  \node[font=\scriptsize,black!65] at ({\xzero+12.5*\bw},2.80) {unused / padding (23 bytes)};

  % Row 2: WirePosted
  \node[font=\footnotesize\ttfamily\bfseries,anchor=east] at (\xzero-0.15,1.80) {WirePosted};
  \draw[thick,sourceblue,fill=sourceblue!18] (\xzero,1.45) rectangle ({\xzero+1*\bw},2.15);
  \node[font=\scriptsize\ttfamily\bfseries,sourceblue!85!black] at ({\xzero+0.5*\bw},1.80) {1};
  \draw[thick,gray!65,fill=gray!12] ({\xzero+1*\bw},1.45) rectangle ({\xzero+8*\bw},2.15);
  \node[font=\scriptsize,black!65] at ({\xzero+4.5*\bw},1.80) {align pad (7 B)};
  \draw[thick,targetgreen!85!black,fill=targetgreen!14] ({\xzero+8*\bw},1.45) rectangle ({\xzero+16*\bw},2.15);
  \node[font=\scriptsize\ttfamily,targetgreen!65!black] at ({\xzero+12.0*\bw},1.80) {tx\_id: u64 (8 B)};
  \draw[thick,quotientorange!90!black,fill=quotientorange!16] ({\xzero+16*\bw},1.45) rectangle ({\xzero+18*\bw},2.15);
  \node[font=\scriptsize\ttfamily,quotientorange!70!black] at ({\xzero+17.0*\bw},1.80) {u16};
  \draw[thick,gray!65,fill=gray!12] ({\xzero+18*\bw},1.45) rectangle ({\xzero+24*\bw},2.15);
  \node[font=\scriptsize,black!65] at ({\xzero+21.0*\bw},1.80) {tail pad (6 B)};

  % Row 3: BatchClosed
  \node[font=\footnotesize\ttfamily\bfseries,anchor=east] at (\xzero-0.15,0.80) {BatchClosed};
  \draw[thick,sourceblue,fill=sourceblue!18] (\xzero,0.45) rectangle ({\xzero+1*\bw},1.15);
  \node[font=\scriptsize\ttfamily\bfseries,sourceblue!85!black] at ({\xzero+0.5*\bw},0.80) {2};
  \draw[thick,gray!65,fill=gray!12] ({\xzero+1*\bw},0.45) rectangle ({\xzero+8*\bw},1.15);
  \node[font=\scriptsize,black!65] at ({\xzero+4.5*\bw},0.80) {align pad (7 B)};
  \draw[thick,targetgreen!85!black,fill=targetgreen!14] ({\xzero+8*\bw},0.45) rectangle ({\xzero+16*\bw},1.15);
  \node[font=\scriptsize\ttfamily,targetgreen!65!black] at ({\xzero+12.0*\bw},0.80) {batch\_id: u64 (8 B)};
  \draw[thick,structurepurple,fill=structurepurple!14] ({\xzero+16*\bw},0.45) rectangle ({\xzero+20*\bw},1.15);
  \node[font=\scriptsize\ttfamily,structurepurple!85!black] at ({\xzero+18.0*\bw},0.80) {count: u32};
  \draw[thick,quotientorange!90!black,fill=quotientorange!16] ({\xzero+20*\bw},0.45) rectangle ({\xzero+24*\bw},1.15);
  \node[font=\scriptsize\ttfamily,quotientorange!70!black] at ({\xzero+22.0*\bw},0.80) {flags: u32};

  \draw[decorate,decoration={brace,mirror,amplitude=5pt},thick,black!75]
    (\xzero,0.32) -- ({\xzero+24*\bw},0.32)
    node[midway,below=6pt,font=\scriptsize]
    {\texttt{size\_of::<LedgerEvent>() == 24} bytes (\texttt{0..=23}), \texttt{align\_of::<LedgerEvent>() == 8}};
\end{tikzpicture}
\caption{Byte layout of \texttt{\#[repr(C, u8)] enum LedgerEvent}. Byte \texttt{0} stores the tag; bytes \texttt{1..=7} align the 8-byte payload region; bytes \texttt{8..=23} overlay variant fields (\texttt{u16} in \texttt{WirePosted} is \texttt{fee\_bps}).}
\label{ch05:fig-enum-byte-layout}
\end{figure}

What if \texttt{BatchClosed} carried \texttt{region: u16} instead of \texttt{flags: u32}, shrinking its payload to $8 + 4 + 2 = 14$ bytes? Under \texttt{\#[repr(C, u8)]}, the enum still takes 24 bytes because the C layout insists on a whole 8-byte-aligned union after the tag. Under default \texttt{repr(Rust)}, \texttt{rustc} is free to lay out each variant around the tag. It keeps the tag at byte~0, slides \texttt{region} into bytes \texttt{2..=3} and \texttt{count} into bytes \texttt{4..=7} (the slots that used to be padding), and puts \texttt{batch\_id} at bytes \texttt{8..=15}. The whole enum fits in \textbf{16 bytes}. Dumping the raw bytes of \texttt{BatchClosed \{ batch\_id: 0x1111..., count: 0x2222\_2222, region: 0x3333 \}} on x86-64 shows exactly that: \texttt{02 00 33 33 22 22 22 22 11 11 11 11 11 11 11 11}.

\subsection{Niche Optimization: Zero-Overhead \texttt{Option<T>}}
\label{ch05:niche}

Why did C and Java put \texttt{NULL} / \texttt{null} inside every pointer type in the first place? Because systems programmers refused to pay 16 bytes (an 8-byte pointer plus a 1-byte tag plus 7 padding bytes) for every nullable pointer in a tree or hash table. Rust achieves \emph{both} static non-null references \emph{and} an 8-byte \texttt{Option<\&T>} through \emph{niche optimization}.

\begin{definition}[Bit-pattern niche]
\label{ch05:def-niche}
Let \texttt{T} occupy $k = \texttt{size\_of::<T>()}$ bytes ($8k$ bits). Any $8k$-bit pattern that is never a valid inhabitant of \texttt{T} under Rust's validity invariants is a \emph{niche} of \texttt{T}.
\end{definition}

Compare \texttt{Option<u64>} against types that possess at least one niche:
\begin{itemize}
  \item \textbf{\texttt{u64} has zero niches.} Every 64-bit pattern from \texttt{0} through \texttt{u64::MAX} is a valid integer. Distinguishing \texttt{None} from \texttt{Some(0)} requires an extra tag byte plus 7 alignment padding bytes, making \texttt{Option<u64>} 16 bytes wide.
  \item \textbf{\texttt{\&T}, \texttt{\&mut T}, \texttt{Box<T>}, \texttt{NonNull<T>}, and \texttt{fn()} forbid \texttt{0x0}.} Because safe references, \texttt{Box}es, and function pointers are never null, the all-zero bit pattern is a niche. Thus \texttt{Option<\&u64>} and \texttt{Option<Box<u64>>} store \texttt{None} as zero and \texttt{Some(p)} as the raw address in \textbf{8 bytes total} --- matching a C pointer at machine level with static null checking on top.
  \item \textbf{\texttt{NonZeroU64} and newtypes.} Ruling out zero with \texttt{core::num::NonZeroU64} creates the same \texttt{0x0} niche for integers (\texttt{size\_of::<Option<NonZeroU64>>() == 8}), and niches propagate through wrapper structs such as \texttt{struct Cents(NonZeroU64)}.
  \item \textbf{\texttt{bool} has 254 niches.} Only \texttt{0x00} and \texttt{0x01} are valid \texttt{bool}s, so \texttt{Option<bool>} encodes \texttt{None} as \texttt{0x02} in \textbf{1 byte}.
\end{itemize}

\begin{lstlisting}[language=Rust]
use core::mem::size_of;
use core::num::NonZeroU64;

const _: () = assert!(size_of::<u64>() == 8);
const _: () = assert!(size_of::<Option<u64>>() == 16);

const _: () = assert!(size_of::<Option<&u64>>() == 8);
const _: () = assert!(size_of::<Option<&mut u64>>() == 8);
const _: () = assert!(size_of::<Option<Box<u64>>>() == 8);
const _: () = assert!(size_of::<Option<NonZeroU64>>() == 8);
const _: () = assert!(size_of::<Option<fn()>>() == 8);
const _: () = assert!(size_of::<Option<bool>>() == 1);
\end{lstlisting}

Figure~\ref{ch05:fig-niche-layout} contrasts the 16-byte layout of \texttt{Option<u64>} with the 8-byte niche layout of \texttt{Option<\&u64>} and \texttt{Option<NonZeroU64>}.

\begin{figure}[H]
\centering
\begin{tikzpicture}[font=\footnotesize,>=Latex]
  % Top panel: Option<u64> = 16 bytes
  \node[font=\bfseries,sourceblue,anchor=west] at (0,3.50) {A. \texttt{Option<u64>} --- 16 bytes (all $2^{64}$ patterns of \texttt{u64} are valid; no niche)};

  \node[font=\scriptsize\ttfamily,anchor=east] at (1.65,2.85) {None};
  \draw[thick,sourceblue,fill=sourceblue!18] (1.80,2.55) rectangle (2.45,3.15);
  \node[font=\scriptsize\ttfamily\bfseries,sourceblue!85!black] at (2.125,2.85) {0};
  \draw[thick,gray!65,fill=gray!12] (2.45,2.55) rectangle (6.95,3.15);
  \node[font=\scriptsize,black!65] at (4.70,2.85) {align pad (1..=7)};
  \draw[thick,gray!65,fill=gray!12] (6.95,2.55) rectangle (12.35,3.15);
  \node[font=\scriptsize,black!65] at (9.65,2.85) {unused payload (8..=15)};

  \node[font=\scriptsize\ttfamily,anchor=east] at (1.65,2.05) {Some(0)};
  \draw[thick,sourceblue,fill=sourceblue!18] (1.80,1.75) rectangle (2.45,2.35);
  \node[font=\scriptsize\ttfamily\bfseries,sourceblue!85!black] at (2.125,2.05) {1};
  \draw[thick,gray!65,fill=gray!12] (2.45,1.75) rectangle (6.95,2.35);
  \node[font=\scriptsize,black!65] at (4.70,2.05) {align pad (1..=7)};
  \draw[thick,targetgreen!85!black,fill=targetgreen!14] (6.95,1.75) rectangle (12.35,2.35);
  \node[font=\scriptsize\ttfamily,targetgreen!65!black] at (9.65,2.05) {0x0000\_0000\_0000\_0000 (8..=15)};

  % Bottom panel: Option<&u64> / Option<NonZeroU64> = 8 bytes
  \node[font=\bfseries,targetgreen!80!black,anchor=west] at (0,1.10) {B. \texttt{Option<\&u64>} or \texttt{Option<NonZeroU64>} --- 8 bytes (niche at \texttt{0x0})};

  \node[font=\scriptsize\ttfamily,anchor=east] at (1.65,0.45) {None};
  \draw[thick,quotientorange!90!black,fill=quotientorange!16] (1.80,0.15) rectangle (7.40,0.75);
  \node[font=\scriptsize\ttfamily,quotientorange!70!black] at (4.60,0.45) {0x0000\_0000\_0000\_0000 (0..=7)};
  \node[font=\scriptsize,black!70,anchor=west] at (7.65,0.45) {forbidden for \texttt{\&u64} / \texttt{NonZeroU64}};

  \node[font=\scriptsize\ttfamily,anchor=east] at (1.65,-0.35) {Some(x)};
  \draw[thick,targetgreen!85!black,fill=targetgreen!14] (1.80,-0.65) rectangle (7.40,-0.05);
  \node[font=\scriptsize\ttfamily,targetgreen!65!black] at (4.60,-0.35) {non-zero ptr or int (0..=7)};
  \node[font=\scriptsize,black!70,anchor=west] at (7.65,-0.35) {zero tag bytes; fits in 1 register};
\end{tikzpicture}
\caption{Tag-beside-payload layout of \texttt{Option<u64>} (16 bytes) versus niche-encoded \texttt{Option<\&u64>} and \texttt{Option<NonZeroU64>} (8 bytes), where \texttt{None} occupies the forbidden all-zero bit pattern.}
\label{ch05:fig-niche-layout}
\end{figure}

\begin{trapbox}
\textbf{The Giant Variant Trap.} Because an \texttt{enum} is sized to its \emph{largest} variant inline, one cold diagnostic variant carrying a 512-byte buffer inflates \emph{every} \texttt{Result} returning that error:
\begin{lstlisting}[language=Rust]
pub enum WireError {
    Timeout,
    InvalidRoutingNumber(u32),
    DiagnosticDump([u8; 512]),        // inline 512-byte array!
}
pub enum SlimWireError {
    Timeout,
    InvalidRoutingNumber(u32),
    DiagnosticDump(Box<[u8; 512]>),   // 8-byte heap pointer on cold path
}
const _: () = assert!(size_of::<Result<u64, WireError>>() == 520);
const _: () = assert!(size_of::<Result<u64, SlimWireError>>() == 16);
\end{lstlisting}
Returning \texttt{Result<u64, WireError>} copies a 520-byte stack slot even when 99.999\% of calls return \texttt{Ok(n)}. Boxing the cold payload shrinks \texttt{Result<u64, SlimWireError>} to 16 bytes --- two CPU registers (\texttt{rax}, \texttt{rdx}) and zero heap allocation on the happy path. Clippy flags oversized error types automatically via \texttt{clippy::result\_large\_err}.
\end{trapbox}

\section{\texttt{Option<T>}, \texttt{Result<T, E>}, and \texttt{?} vs.\ Java~6 \texttt{null} and Checked Exceptions}
\label{ch05:option-result}

Neither \texttt{Option<T>} nor \texttt{Result<T, E>} relies on special compiler privilege for its definition; both are five-line standard-library \texttt{enum}s in \texttt{core}:

\begin{lstlisting}[language=Rust]
pub enum Option<T> {
    None,
    Some(T),
}

#[must_use]
pub enum Result<T, E> {
    Ok(T),
    Err(E),
}
\end{lstlisting}

\subsection{\texttt{Option<T>} and Combinator Pipelines}
\label{ch05:option-combinators}

In Java~6, \texttt{Map.get(id)} returns \texttt{Account} whether the key exists or not; calling a getter such as \texttt{getDailyLimitCents()} on the result compiles without warning, and when the key is absent it fails at runtime with a \texttt{NullPointerException}. In Rust, \texttt{HashMap::get} returns \texttt{Option<\&Account>}, which is a different type from \texttt{\&Account}. You unpack the \texttt{Option} once at the lookup boundary, and downstream functions take plain \texttt{\&Account}.

Rather than nesting four-line \texttt{match} blocks for every optional step, idiomatic Rust chains combinator methods:

\begin{lstlisting}[language=Rust]
use std::collections::HashMap;

pub struct Account {
    pub daily_limit_cents: Option<u64>,
    pub frozen: bool,
}

pub fn effective_limit_cents(
    accounts: &HashMap<String, Account>,
    owner_id: &str,
    default_limit: u64,
) -> u64 {
    accounts
        .get(owner_id)                           // Option<&Account>
        .filter(|acct| !acct.frozen)             // None if account is frozen
        .and_then(|acct| acct.daily_limit_cents) // flatten Option<Option<u64>>
        .map(|cents| cents.min(1_000_000))       // cap at $10,000.00
        .unwrap_or(default_limit)                // fallback u64 on None
}
\end{lstlisting}

Five combinators handle almost every pipeline you will write:
\begin{itemize}
  \item \texttt{.map(f)} transforms \texttt{Some(x)} into \texttt{Some(f(x))} and leaves \texttt{None} untouched.
  \item \texttt{.and\_then(f)} flat-maps when \texttt{f} \emph{itself} returns \texttt{Option<U>}, avoiding nested \texttt{Option<Option<U>>}.
  \item \texttt{.filter(pred)} keeps \texttt{Some(x)} when \texttt{pred(\&x)} is true and turns it into \texttt{None} otherwise.
  \item \texttt{.unwrap\_or(val)} vs.\ \texttt{.unwrap\_or\_else(|| compute())}: function arguments evaluate eagerly, so \texttt{opt.unwrap\_or(load\_default())} executes \texttt{load\_default()} \emph{even when \texttt{opt} is \texttt{Some}}. Use \texttt{unwrap\_or} for cheap constants and \texttt{unwrap\_or\_else} whenever the fallback allocates or computes.
  \item \texttt{.ok\_or(err)} and \texttt{.ok\_or\_else(|| make\_err())} promote \texttt{Option<T>} into \texttt{Result<T, E>}, mapping \texttt{None} to \texttt{Err(err)} so you can propagate it with \texttt{?}.
\end{itemize}

\begin{trapbox}
\textbf{Calling \texttt{.unwrap()} in request handlers re-invites \texttt{NullPointerException}.} Calling \texttt{None.unwrap()} or \texttt{Err(e).unwrap()} triggers a \texttt{panic!}, crashing the current thread or Tokio task. Using \texttt{.unwrap()} on network payloads, query parameters, or database lookups throws away static error handling. Save \texttt{.expect("invariant reason")} for invariants proven true at startup or by an immediately preceding check.
\end{trapbox}

\subsection{\texttt{Result<T, E>} vs.\ C \texttt{errno} and Java~6 Checked Exceptions}
\label{ch05:result-vs-exceptions}

Recoverable error handling in C and Java~6 sits at opposite extremes:
\begin{enumerate}
  \item \textbf{C's \texttt{-1} and \texttt{errno}} return in a single CPU register at zero allocation cost, yet nothing in C's type system warns when a caller ignores the return code and reads an uninitialized output buffer. Rust's \texttt{Result<T, E>} carries \texttt{\#[must\_use]}, so dropping a \texttt{Result} unread triggers an immediate compiler warning.
  \item \textbf{Java~6 checked exceptions (\texttt{throws IOException, SQLException})} put failure in the method signature, yet couple static checking to \emph{stack unwinding} and \emph{stack-trace allocation}. Constructing \texttt{new Exception()} walks every active stack frame inside \texttt{Throwable.fillInStackTrace()}, turning a burst of bad client requests or downstream timeouts into microsecond-scale heap allocation storms.
\end{enumerate}

Rust decouples static error checking from unwinding: \texttt{Result<T, E>} is returned through the exact same register/stack calling convention as any other value, executing a normal \texttt{ret} instruction with zero unwinding tables or heap allocation on the error path.

\subsection{Desugaring the \texttt{?} Operator Step by Step}
\label{ch05:question-mark}

To propagate \texttt{Result} values without repetitive \texttt{if err != nil} blocks, postfix \texttt{?} unwraps \texttt{Ok(v)} inline or returns \texttt{Err} early:

\begin{lstlisting}[language=Rust]
use std::collections::HashMap;
use std::num::ParseIntError;

#[derive(Debug, PartialEq, Eq)]
pub enum ConfigError {
    MissingKey(&'static str),
    InvalidInt(ParseIntError),
    PortZeroOrReserved(u16),
}

impl From<ParseIntError> for ConfigError {
    fn from(err: ParseIntError) -> Self {
        ConfigError::InvalidInt(err)
    }
}

pub fn parse_service_port(
    env_map: &HashMap<String, String>,
) -> Result<u16, ConfigError> {
    let raw = env_map
        .get("LISTEN_PORT")
        .ok_or(ConfigError::MissingKey("LISTEN_PORT"))?;
    let port: u16 = raw.trim().parse()?;
    if port < 1024 {
        return Err(ConfigError::PortZeroOrReserved(port));
    }
    Ok(port)
}
\end{lstlisting}

Notice how \texttt{raw.trim().parse()?} produces a \texttt{ParseIntError}, yet returns cleanly from a function returning \texttt{Result<u16, ConfigError>}. Ignoring trait plumbing (\texttt{core::ops::Try}), \texttt{expr?} on a \texttt{Result} expands to:

\begin{lstlisting}[language=Rust]
match expr {
    Ok(val) => val,
    Err(err) => return Err(core::convert::From::from(err)),
}
\end{lstlisting}

Because we implemented \texttt{From<ParseIntError> for ConfigError}, \texttt{From::from(err)} wraps the low-level parse error into \texttt{ConfigError::InvalidInt(err)} automatically; when the error already has type \texttt{ConfigError}, the standard identity \texttt{impl<T> From<T> for T} compiles to a no-op. Figure~\ref{ch05:fig-question-mark-flow} traces both value tracks.

\begin{figure}[H]
\centering
\begin{tikzpicture}[font=\footnotesize,>=Latex]
  % Top track: 4 boxes with generous 1.15 cm gaps (total width ~14.2 cm)
  \node[draw,thick,sourceblue,fill=sourceblue!8,rounded corners=3pt,minimum width=2.25cm,minimum height=0.85cm,align=center] (get) at (1.15,2.10)
    {\texttt{env\_map.get}\\$\to$ \texttt{Option}};
  \node[draw,thick,sourceblue,fill=sourceblue!8,rounded corners=3pt,minimum width=2.35cm,minimum height=0.85cm,align=center] (okor) at (4.75,2.10)
    {\texttt{.ok\_or(...)}\\$\to$ \texttt{Result}};
  \node[draw,thick,sourceblue,fill=sourceblue!8,rounded corners=3pt,minimum width=2.55cm,minimum height=0.85cm,align=center] (parse) at (8.55,2.10)
    {\texttt{raw.parse::<u16>()}\\$\to$ \texttt{Result}};
  \node[draw,thick,targetgreen!85!black,fill=targetgreen!12,rounded corners=3pt,minimum width=2.35cm,minimum height=0.85cm,align=center] (done) at (12.35,2.10)
    {\texttt{port >= 1024}\\$\to$ \texttt{Ok(port)}};

  % Bottom early-return collector bar spanning beneath okor and parse
  \node[draw,thick,trapred,fill=trapred!10,rounded corners=3pt,minimum width=7.2cm,minimum height=0.78cm,align=center] (errbar) at (6.85,0.30)
    {Early return \texttt{Err(ConfigError)} in registers (zero stack unwinding)};

  % Ok track horizontal arrows
  \draw[->,thick,targetgreen!80!black] (get) -- node[above=2pt,font=\scriptsize\ttfamily] {Some} (okor);
  \draw[->,thick,targetgreen!80!black] (okor) -- node[above=2pt,font=\scriptsize\ttfamily] {Ok(raw)} (parse);
  \draw[->,thick,targetgreen!80!black] (parse) -- node[above=2pt,font=\scriptsize\ttfamily] {Ok(port)} (done);

  % Err drop arrows straight down into errbar
  \draw[->,thick,trapred] (okor.south) -- (4.75,0.69)
    node[midway,left=3pt,font=\scriptsize\ttfamily,trapred] {? on Err};
  \draw[->,thick,quotientorange!90!black] (parse.south) -- (8.55,0.69)
    node[midway,right=3pt,font=\scriptsize\ttfamily,quotientorange!75!black] {? via From::from};
\end{tikzpicture}
\caption{Two-track control flow of \texttt{?} inside \texttt{parse\_service\_port}. Each \texttt{?} unwraps \texttt{Ok} rightward or converts \texttt{Err} via \texttt{From::from} into a normal register return.}
\label{ch05:fig-question-mark-flow}
\end{figure}

The \texttt{?} operator also works on \texttt{Option<T>} inside functions returning \texttt{Option<U>}:

\begin{lstlisting}[language=Rust]
pub fn header_fee_bps(lines: &[&str]) -> Option<u16> {
    let first = lines.first()?;                  // None -> early return None
    let digits = first.strip_prefix("FEE_BPS=")?; // None -> early return None
    digits.trim().parse::<u16>().ok()            // Result<u16, _> -> Option<u16>
}
\end{lstlisting}

Rust intentionally rejects mixing \texttt{?} on \texttt{Option} directly inside a \texttt{Result}-returning function because \texttt{None} carries no error value \texttt{E}; bridge explicitly with \texttt{.ok\_or(...)} or \texttt{.ok()}.

\section{Exhaustive Pattern Matching as Compiler-Enforced Case Analysis}
\label{ch05:pattern-matching}

Pattern matching is the exact dual of data construction: it inspects a composite value's shape and binds variables to its interior pieces in one checked step.

\subsection{\texttt{match} Expressions, Destructuring, \texttt{@} Bindings, and Guards}
\label{ch05:match-anatomy}

Unlike C and Java~6 \texttt{switch} statements --- limited to primitive integers or enum constants, prone to accidental fallthrough without \texttt{break;}, and unable to produce a value --- Rust's \texttt{match} is an expression with zero fallthrough that destructures arbitrary nested shapes:

\begin{lstlisting}[language=Rust]
#[derive(Debug, Clone)]
pub struct WireFrame {
    pub channel: u8,
    pub sender: String,
    pub payload: Vec<u8>,
}

pub fn classify_frame(frame: &WireFrame, fast_limit: usize) -> String {
    match frame {
        WireFrame {
            channel: ch @ 1..=9,
            sender,
            payload,
        } if payload.len() <= fast_limit => match payload.as_slice() {
            [] => format!("priority ch {ch} ({sender}): keepalive"),
            [opcode] => format!("priority ch {ch} ({sender}): ctrl 0x{opcode:02x}"),
            [opcode, body @ ..] => format!(
                "priority ch {ch} ({sender}): op 0x{opcode:02x}, {}B",
                body.len()
            ),
        },
        WireFrame { channel: 0, sender, .. } => format!("sys ({sender})"),
        WireFrame { channel, sender, payload } => {
            format!("std ch {channel} ({sender}): {}B", payload.len())
        }
    }
}
\end{lstlisting}

Four pattern tools work together in \texttt{classify\_frame}:
\begin{enumerate}
  \item \textbf{Ranges and \texttt{@} bindings (\texttt{channel: ch @ 1..=9}):} \texttt{1..=9} tests inclusive range membership while \texttt{ch @} binds the matched byte so the arm body can print \texttt{ch}.
  \item \textbf{Exhaustive slice patterns (\texttt{[]}, \texttt{[opcode]}, \texttt{[opcode, body @ ..]}):} destructure slices of length $0$, $1$, and $\ge 2$ without manual index bounds checks.
  \item \textbf{Rest patterns (\texttt{..}):} ignore remaining struct fields (\texttt{WireFrame \{ channel: 0, sender, .. \}}) or bind remaining slice elements (\texttt{body @ ..} as \texttt{\&[u8]}).
  \item \textbf{Match guards (\texttt{if payload.len() <= fast\_limit}):} attach a runtime boolean filter. Because the compiler cannot prove an arbitrary runtime boolean is always true, a guarded arm never satisfies exhaustiveness by itself; frames failing the guard continue to later arms.
\end{enumerate}

\subsection{Match Ergonomics: Borrowing Without Moving}
\label{ch05:match-ergonomics}

Trace ownership in \texttt{classify\_frame} using Chapters~3 and~4: \texttt{frame} has type \texttt{\&WireFrame}, and \texttt{sender: String} is not \texttt{Copy}. Why didn't \texttt{WireFrame \{ sender, .. \}} fail to move \texttt{sender} out of a shared reference? Under \emph{match ergonomics}, matching a value pattern against \texttt{\&T} (or \texttt{\&mut T}) pushes the reference inward so every variable binds as \texttt{\&Field} (or \texttt{\&mut Field}) automatically: \texttt{sender} gets type \texttt{\&String} and \texttt{payload} gets \texttt{\&Vec<u8>}. You only need explicit \texttt{ref} or \texttt{ref mut} when matching on an \emph{owned} value where you wish to borrow one field without consuming the rest (e.g., \texttt{Some(ref s)} on an owned \texttt{Option<String>}).

\subsection{Exhaustiveness Across Codebase Evolution and Control-Flow Shortcuts}
\label{ch05:exhaustiveness}

Suppose \texttt{enum OrderStatus \{ Pending, Settled, Failed \}} is matched in forty files across a service. Six months later, a teammate adds \texttt{Disputed \{ dispute\_id: u64 \}}. In Java~6, existing \texttt{switch (status)} blocks with a \texttt{default:} branch compile cleanly and silently misroute disputed orders in production. In Rust, \texttt{cargo check} immediately flags \texttt{non-exhaustive patterns: `OrderStatus::Disputed \{ .. \}` not covered} at every single unhandled \texttt{match} across all forty files, which gives you a list of every place that must decide how to handle the new state.

\begin{trapbox}
\textbf{Avoid wildcard \texttt{\_ => ...} arms on your own domain enums.} A catch-all \texttt{\_} arm disables future exhaustiveness errors on that \texttt{match}. Group shared cases with an explicit or-pattern instead (\texttt{OrderStatus::Settled | OrderStatus::Failed => false}), reserving \texttt{\_} for scalars (\texttt{u8}, \texttt{\&str}) or external \texttt{\#[non\_exhaustive]} enums.
\end{trapbox}

When a full multi-arm \texttt{match} is more ceremony than a single check needs, reach for \texttt{let-else} guard clauses or \texttt{matches!}:

\begin{lstlisting}[language=Rust]
pub fn parse_bearer_token(header: Option<&str>) -> Result<&str, &'static str> {
    let Some(raw) = header else {
        return Err("missing Authorization header");
    };
    let Some(token) = raw.strip_prefix("Bearer ") else {
        return Err("Authorization scheme must be Bearer");
    };
    let token = token.trim();
    if token.is_empty() {
        return Err("empty bearer token");
    }
    Ok(token)
}

pub fn is_retryable_http(code: Option<u16>) -> bool {
    matches!(code, Some(429 | 500..=599))
}
\end{lstlisting}

In \texttt{let Some(raw) = header else \{ return ...; \};}, \texttt{raw} stays bound at the outer indentation level with zero rightward drift because the compiler verifies that the \texttt{else} block diverges to type \texttt{!}.

\begin{focusbox}
\textbf{Choosing the right pattern construct:}
\begin{itemize}
  \item \textbf{\texttt{match expr \{ ... \}}}: multi-variant dispatch and domain enums where future variant additions must trigger compiler errors.
  \item \textbf{\texttt{let Pat = expr else \{ return ...; \};}}: flat guard clauses unpacking the happy path into outer scope.
  \item \textbf{\texttt{if let Pat = expr \{ ... \}}}: optional side effect on a single pattern when all other cases truly do nothing.
  \item \textbf{\texttt{matches!(expr, Pat)}}: concise boolean test supporting ranges, destructuring, and \texttt{|} alternatives.
\end{itemize}
\end{focusbox}

\section{Case Study: Refactoring a Legacy Java~6 \texttt{PaymentOrder} into a Rust State Machine}
\label{ch05:case-study}

Let us combine cardinality algebra, niche newtypes, \texttt{Result}, and exhaustive \texttt{match} to refactor a buggy Java~6 payment order class into a compile-time-checked Rust state machine.

\subsection{Four Silent Bugs in the Legacy Java~6 \texttt{PaymentOrder}}
\label{ch05:java-payment-bugs}

\begin{lstlisting}[language=Java]
// Java 6: Mutable product class with nullable lifecycle fields
public class PaymentOrder {
    public enum Status { CREATED, AUTHORIZED, CAPTURED, SETTLED, REFUNDED, FAILED }

    private final long orderId;
    private Status status;
    private long requestedCents, authorizedCents, capturedCents, settledCents;
    private String authCode;        // null in CREATED, FAILED
    private String captureTxId;     // null until CAPTURED
    private Long settledAtEpochMs;  // null until SETTLED
    private String refundReason;    // null unless REFUNDED
    private String failureCode;     // null unless FAILED

    public PaymentOrder(long orderId, long requestedCents) {
        this.orderId = orderId;
        this.status = Status.CREATED;
        this.requestedCents = requestedCents;
    }

    public void capture(long amountCents, String txId) {
        this.capturedCents = amountCents;
        this.status = Status.CAPTURED;
        // Bug 1: forgot this.captureTxId = txId -> NPE in nightly settlement!
        // Bug 2: forgot status == AUTHORIZED and amountCents <= authorizedCents!
    }

    public void refund(long amountCents, String reason) {
        if (amountCents > this.settledCents) {
            throw new IllegalArgumentException("refund exceeds settled");
        }
        this.refundReason = reason;
        this.status = Status.REFUNDED; // Bug 3: refund(0, "x") succeeds on CREATED!
    }
}
\end{lstlisting}

Java's compiler accepts all four failure modes here without complaint: (1)~moving to \texttt{CAPTURED} while leaving \texttt{captureTxId == null}, which surfaces hours later as a \texttt{NullPointerException} in the nightly settlement job; (2)~capturing an order that was never authorized, or capturing more than was authorized, because \texttt{capture} checks neither \texttt{status} nor \texttt{authorizedCents}; (3)~calling \texttt{refund(0, "x")} on an uncharged \texttt{CREATED} order, which passes the guard because \texttt{settledCents} defaulted to \texttt{0L}; and (4)~accepting zero or negative amounts, since every amount is a signed \texttt{long}. On top of that, a method that throws halfway through leaves a half-mutated object behind.

\subsection{Step-by-Step Refactoring into an Algebraic State Machine}
\label{ch05:rust-payment-machine}

\textbf{Step 1: Validate leaf primitives with niche-preserving newtypes.} Wrapping \texttt{NonZeroU64} in \texttt{Cents} rules out zero and negative amounts at construction while keeping \texttt{Option<Cents>} at 8 bytes:

\begin{lstlisting}[language=Rust]
use core::num::NonZeroU64;

#[derive(Debug, Clone, Copy, PartialEq, Eq)]
pub struct Cents(NonZeroU64);

impl Cents {
    pub fn new(raw: u64) -> Option<Self> { NonZeroU64::new(raw).map(Cents) }
    pub fn get(self) -> u64 { self.0.get() }
}
const _: () = assert!(core::mem::size_of::<Option<Cents>>() == 8);

#[derive(Debug, Clone, PartialEq, Eq)]
pub struct AuthCode(String);
impl AuthCode {
    pub fn parse(raw: &str) -> Option<Self> {
        let s = raw.trim();
        (!s.is_empty()).then(|| AuthCode(s.to_owned()))
    }
}

#[derive(Debug, Clone, PartialEq, Eq)]
pub struct CaptureTxId(String);
impl CaptureTxId {
    pub fn parse(raw: &str) -> Option<Self> {
        let s = raw.trim();
        (!s.is_empty()).then(|| CaptureTxId(s.to_owned()))
    }
}
\end{lstlisting}

\textbf{Step 2: Model the lifecycle as a sum of exact per-stage payloads.} Each variant of \texttt{PaymentState} carries \emph{only} the data that exists in that stage, matching Figure~\ref{ch05:fig-payment-state-machine}:

\begin{figure}[H]
\centering
\begin{tikzpicture}[font=\scriptsize,>=Latex]
  % Two-row spacious layout (total width ~13.8 cm, generous 1.7 cm arrows!)
  \node[draw,thick,sourceblue,fill=sourceblue!8,rounded corners=3pt,minimum width=2.4cm,minimum height=1.1cm,align=center] (created) at (1.5,2.35)
    {\textbf{\texttt{Created}}\\ \texttt{requested}};

  \node[draw,thick,sourceblue,fill=sourceblue!8,rounded corners=3pt,minimum width=2.6cm,minimum height=1.1cm,align=center] (auth) at (5.9,2.35)
    {\textbf{\texttt{Authorized}}\\ \texttt{authorized, auth\_code}};

  \node[draw,thick,sourceblue,fill=sourceblue!8,rounded corners=3pt,minimum width=2.8cm,minimum height=1.1cm,align=center] (cap) at (10.6,2.35)
    {\textbf{\texttt{Captured}}\\ \texttt{captured, capture\_tx\_id}};

  % Bottom row: Failed on left, Settled under Captured, Refunded on right
  \node[draw,thick,trapred,fill=trapred!10,rounded corners=3pt,minimum width=2.8cm,minimum height=1.0cm,align=center] (fail) at (3.7,0.35)
    {\textbf{\texttt{Failed}}\\ \texttt{failure\_code}};

  \node[draw,thick,targetgreen!85!black,fill=targetgreen!10,rounded corners=3pt,minimum width=2.8cm,minimum height=1.0cm,align=center] (set) at (8.5,0.35)
    {\textbf{\texttt{Settled}}\\ \texttt{settled, settled\_at\_ms}};

  \node[draw,thick,structurepurple,fill=structurepurple!10,rounded corners=3pt,minimum width=2.6cm,minimum height=1.0cm,align=center] (ref) at (13.0,0.35)
    {\textbf{\texttt{Refunded}}\\ \texttt{refunded, reason}};

  % Transitions with generous arrow clearance
  \draw[->,thick,sourceblue!90!black] (created) -- node[above=2pt,font=\scriptsize\ttfamily] {authorize} (auth);
  \draw[->,thick,sourceblue!90!black] (auth) -- node[above=2pt,font=\scriptsize\ttfamily] {capture} (cap);
  \draw[->,thick,targetgreen!80!black] (cap.south) -- (8.9,0.85) node[midway,right=4pt,font=\scriptsize\ttfamily] {settle};
  \draw[->,thick,structurepurple] (set) -- node[above=2pt,font=\scriptsize\ttfamily] {refund} (ref);

  % Failure transitions ONLY from Created and Authorized
  \draw[->,thick,trapred] (created.south) -- (2.9,0.85) node[midway,left=3pt,font=\scriptsize\ttfamily] {fail};
  \draw[->,thick,trapred] (auth.south) -- (4.5,0.85) node[midway,right=3pt,font=\scriptsize\ttfamily] {fail};
\end{tikzpicture}
\caption{Lifecycle graph of \texttt{PaymentOrder}. Each variant owns only the fields valid in that state, and \texttt{fail} is reachable only before funds are captured.}
\label{ch05:fig-payment-state-machine}
\end{figure}

\textbf{Step 3: Implement ownership-consuming transitions with explicit exhaustive arms.} Bringing variants into scope with \texttt{use PaymentState::*;} keeps each \texttt{match} arm concise while spelling out rejected states via or-patterns (rather than \texttt{\_}) so adding a new state later triggers compiler errors across every transition:

\begin{lstlisting}[language=Rust]
#[derive(Debug, Clone, PartialEq, Eq)]
pub enum PaymentState {
    Created { requested: Cents },
    Authorized { authorized: Cents, auth_code: AuthCode },
    Captured { authorized: Cents, captured: Cents, capture_tx_id: CaptureTxId },
    Settled { settled: Cents, capture_tx_id: CaptureTxId, settled_at_ms: u64 },
    Refunded { refunded: Cents, capture_tx_id: CaptureTxId, reason: String },
    Failed { failure_code: String },
}

#[derive(Debug, Clone, PartialEq, Eq)]
pub struct PaymentOrder {
    pub order_id: u64,
    pub state: PaymentState,
}

#[derive(Debug, Clone, PartialEq, Eq)]
pub enum OrderTransitionError {
    InvalidState { attempted: &'static str },
    CaptureExceedsAuth { authorized: u64, requested: u64 },
    RefundExceedsSettled { settled: u64, requested: u64 },
    EmptyReason,
}

use PaymentState::*;

impl PaymentOrder {
    pub fn new(order_id: u64, requested: Cents) -> Self {
        Self { order_id, state: Created { requested } }
    }

    pub fn authorize(self, auth_code: AuthCode) -> Result<Self, OrderTransitionError> {
        match self.state {
            Created { requested } => Ok(Self {
                order_id: self.order_id,
                state: Authorized { authorized: requested, auth_code },
            }),
            Authorized { .. } | Captured { .. } | Settled { .. }
            | Refunded { .. } | Failed { .. } => {
                Err(OrderTransitionError::InvalidState { attempted: "authorize" })
            }
        }
    }

    pub fn capture(
        self,
        amount: Cents,
        capture_tx_id: CaptureTxId,
    ) -> Result<Self, OrderTransitionError> {
        match self.state {
            Authorized { authorized, .. } => {
                if amount.get() > authorized.get() {
                    return Err(OrderTransitionError::CaptureExceedsAuth {
                        authorized: authorized.get(),
                        requested: amount.get(),
                    });
                }
                Ok(Self {
                    order_id: self.order_id,
                    state: Captured { authorized, captured: amount, capture_tx_id },
                })
            }
            Created { .. } | Captured { .. } | Settled { .. }
            | Refunded { .. } | Failed { .. } => {
                Err(OrderTransitionError::InvalidState { attempted: "capture" })
            }
        }
    }

    pub fn settle(self, settled_at_ms: u64) -> Result<Self, OrderTransitionError> {
        match self.state {
            Captured { captured, capture_tx_id, .. } => Ok(Self {
                order_id: self.order_id,
                state: Settled { settled: captured, capture_tx_id, settled_at_ms },
            }),
            Created { .. } | Authorized { .. } | Settled { .. }
            | Refunded { .. } | Failed { .. } => {
                Err(OrderTransitionError::InvalidState { attempted: "settle" })
            }
        }
    }

    pub fn refund(self, amount: Cents, reason: String) -> Result<Self, OrderTransitionError> {
        if reason.trim().is_empty() {
            return Err(OrderTransitionError::EmptyReason);
        }
        match self.state {
            Settled { settled, capture_tx_id, .. } => {
                if amount.get() > settled.get() {
                    return Err(OrderTransitionError::RefundExceedsSettled {
                        settled: settled.get(),
                        requested: amount.get(),
                    });
                }
                Ok(Self {
                    order_id: self.order_id,
                    state: Refunded { refunded: amount, capture_tx_id, reason },
                })
            }
            Created { .. } | Authorized { .. } | Captured { .. }
            | Refunded { .. } | Failed { .. } => {
                Err(OrderTransitionError::InvalidState { attempted: "refund" })
            }
        }
    }

    pub fn fail(self, failure_code: String) -> Result<Self, OrderTransitionError> {
        match self.state {
            Created { .. } | Authorized { .. } => Ok(Self {
                order_id: self.order_id,
                state: Failed { failure_code },
            }),
            Captured { .. } | Settled { .. } | Refunded { .. } | Failed { .. } => {
                Err(OrderTransitionError::InvalidState { attempted: "fail" })
            }
        }
    }
}
\end{lstlisting}

Each of the four Java~6 bugs is now either unrepresentable or rejected by a checked transition: \texttt{Captured} cannot be constructed without a non-empty \texttt{CaptureTxId}; \texttt{capture} only accepts an \texttt{Authorized} order and returns \texttt{CaptureExceedsAuth} instead of over-capturing; \texttt{refund} can only read a balance out of \texttt{PaymentState::Settled}; and \texttt{Cents} cannot hold zero or a negative number. Because each transition consumes \texttt{self} and builds a fresh value, there is no half-mutated object to leave behind, and callers cannot keep using a stale pre-transition variable. Moving \texttt{capture\_tx\_id} from \texttt{Captured} into \texttt{Settled} copies 24 stack bytes without cloning the heap characters.

\begin{connectionbox}
\textbf{Recovering \texttt{self} on error, and compile-time typestate.} If a caller needs the un-transitioned \texttt{PaymentOrder} back when a transition fails without cloning up front, return \texttt{Result<Self, (Self, OrderTransitionError)>}. The standard library uses the same ownership-recovery idea: \texttt{String::from\_utf8} hands the bytes back inside its \texttt{FromUtf8Error}, and \texttt{mpsc::Sender::send} hands the unsent value back inside its \texttt{SendError}. In Chapter~8 we will also meet the \emph{typestate} pattern (\texttt{PaymentOrder<Created>}, \texttt{PaymentOrder<Authorized>}), which encodes known states in generic type parameters; runtime sum enums (\texttt{PaymentState}, essential when loading rows dynamically from PostgreSQL in Chapter~13) and compile-time typestates work side by side in production Rust.
\end{connectionbox}

In Chapter~6 we move from modeling the \emph{shape} of a single type to abstracting \emph{shared behavior} across types: traits, generics, monomorphization, and \texttt{dyn Trait} vtables --- the machinery behind \texttt{From::from} and \texttt{?} in this chapter, built without a single \texttt{extends} superclass.

\section{Exercises}
\label{ch05:exercises}

\begin{exercise}[Counting inhabitants vs.\ domain states]
\label{ch05:ex-cardinality}
A Java~6 \texttt{Subscription} class stores \texttt{boolean isTrial; boolean isCanceled; Integer cancelCode;}, whereas in the domain a subscription is either in a free trial, active paid, or canceled with a 16-bit code (\texttt{0..=65535}).
(a) Compute the exact representable cardinality of the Java~6 class (remembering \texttt{null} inside \texttt{Integer}).
(b) Compute the exact cardinality of valid domain states.
(c) Write a Rust \texttt{enum SubscriptionState} whose cardinality equals~(b) \emph{exactly}.
\end{exercise}

\begin{exercise}[Algebraic identities with \texttt{!} and \texttt{()}]
\label{ch05:ex-never-algebra}
Compute $\abs{\texttt{T}}$ for (a)~\texttt{Result<bool, !>}, (b)~\texttt{(u8, ())}, (c)~\texttt{(u64, !)}, and (d)~\texttt{Option<Option<bool>>}, giving a simpler isomorphic Rust type for each. Then explain why \texttt{fn witness(x: (u64, !)) -> String \{ match x.1 \{\} \}} compiles with zero \texttt{match} arms and can never be called at runtime.
\end{exercise}

\begin{exercise}[Struct padding by hand and with \texttt{rustc}]
\label{ch05:ex-struct-padding}
Given fields \texttt{a: u8, b: u32, c: u16, d: u64} on a 64-bit target:
(a) Under \texttt{\#[repr(C)]} in order \texttt{(a, b, c, d)}, list every field offset, padding range, \texttt{align\_of}, and \texttt{size\_of}.
(b) Reorder the fields manually under \texttt{\#[repr(C)]} to minimize \texttt{size\_of}, and verify both structs with compile-time \texttt{assert!}s in \texttt{rustc}.
\end{exercise}

\begin{exercise}[Tagged union byte layout and tag tucking]
\label{ch05:ex-tagged-union}
Define \texttt{MarketTick} with three variants: \texttt{Halted}, \texttt{Quote \{ symbol\_id: u32, bid: u64, ask: u64 \}}, and \texttt{Trade \{ symbol\_id: u32, price: u64, lots: u16 \}}.
(a) Find \texttt{align\_of} and \texttt{size\_of} under \texttt{\#[repr(C, u8)]} and give every field's byte offset.
(b) Under default \texttt{repr(Rust)}, \texttt{MarketTick} shrinks from 32 bytes to 24 bytes. Explain where inside those 24 bytes \texttt{rustc} tucks the discriminant tag, and verify with \texttt{rustc}.
\end{exercise}

\begin{exercise}[Hunting for niches]
\label{ch05:ex-niches}
Predict \texttt{size\_of} on 64-bit Rust before checking with \texttt{rustc}, justifying each answer via niches:
(a)~\texttt{Option<Option<bool>>};
(b)~\texttt{Option<\&[u8]>} (recall from Chapter~3 that \texttt{\&[u8]} is a 16-byte \texttt{(ptr, len)} pair);
(c)~\texttt{Option<Option<NonZeroU64>>};
(d)~\texttt{Result<\&u64, ()>}.
Explain why (c) grows to 16 bytes while (a) stays 1 byte and (d) stays 8 bytes.
\end{exercise}

\begin{exercise}[Diagnosing and curing a Giant Variant Trap]
\label{ch05:ex-giant-variant}
A middleware returns \texttt{Result<UserId, AuthError>}, where \texttt{UserId} wraps a \texttt{NonZeroU64}. The \texttt{AuthError} variants are \texttt{MissingCookie}, \texttt{Expired \{ at\_s: u64 \}}, and \texttt{InvalidJwt \{ raw: [u8; 1024], pos: usize \}}.
(a) Find the size of \texttt{Result<UserId, AuthError>} and explain why it slows down the 99.99\% \texttt{Ok} fast path.
(b) Refactor \texttt{AuthError} so the \texttt{Result} drops to 16 bytes while keeping all diagnostic data, and prove it with a \texttt{const} assertion.
\end{exercise}

\begin{exercise}[Flattening nested matches with \texttt{?}, \texttt{From}, and \texttt{let-else}]
\label{ch05:ex-desugar-refactor}
Write a function \texttt{parse\_\allowbreak{}ledger\_\allowbreak{}line} that takes \texttt{raw: Option<\&str>} and returns a \texttt{Result} whose success value is the pair \texttt{(NonZeroU64, i64)} and whose error type is your own \texttt{LineParseError}. Reject \texttt{None}, split once on \texttt{':'}, parse the left side as \texttt{NonZeroU64} and the right as \texttt{i64}, and use \texttt{let-else}, \texttt{?}, and \texttt{From} with \emph{zero} nested \texttt{match} blocks. Write out the manual \texttt{match} expansion of one \texttt{?} in your solution.
\end{exercise}

\begin{exercise}[Match ergonomics and accidental moves]
\label{ch05:ex-match-ergonomics}
Let \texttt{JobOutcome} have variants \texttt{Done(u64)} and \texttt{Aborted(String)}, and let \texttt{out} have type \texttt{\&mut JobOutcome}. Explain why matching on \texttt{*out} with arm \texttt{JobOutcome::Aborted(r) => r.len()} fails to compile whereas matching on \texttt{out} directly succeeds. Give the exact type of \texttt{r} in the working code and two ways to obtain \texttt{\&String} instead.
\end{exercise}

\begin{exercise}[Extending the \texttt{PaymentOrder} state machine]
\label{ch05:ex-extend-payment-order}
Extend \texttt{PaymentState} from Section~\ref{ch05:case-study} with two new variants: \texttt{PartiallyRefunded} (carrying \texttt{settled}, \texttt{refunded\_so\_far}, and \texttt{capture\_tx\_id} while $0 < R < S$, moving to \texttt{Refunded} once $R == S$) and \texttt{Disputed} (carrying \texttt{dispute\_id: u64} and \texttt{capture\_tx\_id}, entered from \texttt{Settled} or \texttt{PartiallyRefunded} via \texttt{open\_dispute} and blocking manual refunds). Compile before editing any methods to watch exhaustiveness flag all five transitions, then implement \texttt{refund} and \texttt{open\_dispute} with unit assertions.
\end{exercise}
